% Generated draft: requires mathematical and source review.
\documentclass[11pt]{article}
\usepackage[margin=25mm]{geometry}
\usepackage{fontspec}
\setmainfont{Latin Modern Roman}
\usepackage{amsmath,amssymb,mathtools,unicode-math}
\setmathfont{Latin Modern Math}
\usepackage{enumitem,needspace,xurl,hyperref}
\setlength{\parindent}{0pt}
\setlength{\parskip}{5pt}
\setlength{\emergencystretch}{3em}
\newcommand{\N}{\mathbb{N}}
\newcommand{\Z}{\mathbb{Z}}
\newcommand{\Q}{\mathbb{Q}}
\newcommand{\R}{\mathbb{R}}
\newcommand{\C}{\mathbb{C}}
\newcommand{\F}{\mathbb{F}}
\newcommand{\A}{\mathbb{A}}
\newcommand{\PP}{\mathbb{P}}
\newcommand{\E}{\mathbb{E}}
\newcommand{\eps}{\varepsilon}
\newcommand{\dd}{\,\mathrm{d}}
\begin{document}
\section*{UnsolvedMath ID: 700838}
\textbf{Source problem code:} catalogue entry
\subsection{Definitions and mathematical statement}
For each positive integer \(r\), let \(F_r\) denote a free group of rank \(r\). An epimorphism here means a surjective group homomorphism, and \(\operatorname{Aut}(F_r)\) denotes the group of automorphisms of \(F_r\).

Is the following statement true? For every integer \(n\geq 2\) and every pair of epimorphisms
\[
\varphi,\psi:F_{2n}\twoheadrightarrow F_n\times F_n,
\]
there exists \(\alpha\in\operatorname{Aut}(F_{2n})\) such that
\[
\varphi\circ\alpha=\psi,
\qquad\text{that is,}\qquad
\varphi(\alpha(g))=\psi(g)\quad\text{for every }g\in F_{2n}.
\]
The source writes this relation as \(\alpha\varphi=\psi\); the displayed formulation makes explicit that the automorphism acts on the domain before the epimorphism. The target group is fixed, and the statement allows no additional automorphism of the target.
\subsection{Short English statement}
For every integer \(n\geq 2\), can any two surjective homomorphisms from the free group of rank \(2n\) onto \(F_n\times F_n\) be related by precomposing one with an automorphism of the domain?
\subsection{Sources}
\begin{itemize}
\item Local UnsolvedMath registry, supplied complete source record: permanent ID \(700838\), exact source problem code \texttt{catalogue entry}, title \emph{Kourovka Notebook Problem 12.19}. The supplied registry status is \texttt{open}; this status has not been independently verified.
\item \emph{The Kourovka Notebook (Unsolved Problems in Group Theory)}, original problem 12.19, 12th issue (1992); the supplied background also identifies the 21st-issue edition (2026), arXiv identifier \texttt{1401.0300v48}. The supplied collection description names E. I. Khukhro and V. D. Mazurov as editors.
\item Proposed by R. I. Grigorchuk, with proposed year 1992 according to the supplied record.
\end{itemize}
\end{document}
